% Created 2026-04-23 Thu 10:57
% Intended LaTeX compiler: pdflatex
\documentclass{beamer}\usepackage{listings}
\usepackage{color}
\usepackage{amsmath}
\usepackage{array}
\usepackage[T1]{fontenc}
\usepackage{natbib}
\lstset{
keywordstyle=\color{blue},
commentstyle=\color{red},stringstyle=\color[rgb]{0,.5,0},
basicstyle=\ttfamily\small,
columns=fullflexible,
breaklines=true,
breakatwhitespace=false,
numbers=left,
numberstyle=\ttfamily\tiny\color{gray},
stepnumber=1,
numbersep=10pt,
backgroundcolor=\color{white},
tabsize=4,
literate={~}{$\sim$}{1}{ø}{{\o}}1{æ}{{\ae}}1{å}{{\aa}}1{Ø}{{\OE}}1{Æ}{{\AE}}1{Å}{{\AA}}1,keepspaces=true,
showspaces=false,
showstringspaces=false,
xleftmargin=.23in,
frame=single,
basewidth={0.5em,0.4em},
}
\institute{Steno Diabetes Center Copenhagen}
\subtitle{}
\RequirePackage[absolute, overlay]{textpos}
\setbeamertemplate{footline}[frame number]
\setbeamertemplate{navigation symbols}{}
\RequirePackage{fancyvrb}
\RequirePackage{array}
\RequirePackage{multirow}
\RequirePackage{tcolorbox}
\definecolor{mygray}{rgb}{.95, 0.95, 0.95}
\newcommand{\mybox}[1]{\vspace{.5em}\begin{tcolorbox}[boxrule=0pt,colback=mygray] #1 \end{tcolorbox}}
\newcommand{\sfootnote}[1]{\renewcommand{\thefootnote}{\fnsymbol{footnote}}\footnote{#1}\setcounter{footnote}{0}\renewcommand{\thefootnote}{\arabic{foot note}}}
\usepackage{alphalph}
\alphalph{\value{footnote}}
\setbeamertemplate{itemize item}{\textbullet}
\setbeamertemplate{itemize subitem}{-}
\setbeamertemplate{itemize subsubitem}{}
\setbeamertemplate{enumerate item}{\insertenumlabel.}
\setbeamertemplate{enumerate subitem}{\insertenumlabel.\insertsubenumlabel}
\setbeamertemplate{enumerate subsubitem}{\insertenumlabel.\insertsubenumlabel.\insertsubsubenumlabel}
\setbeamertemplate{enumerate mini template}{\insertenumlabel}
\makeatletter\def\blfootnote{\xdef\@thefnmark{}\@footnotetext}\makeatother
\newcommand{\E}{\ensuremath{\mathrm{E}}}
\renewcommand{\P}{\ensuremath{\mathrm{P}}}
\renewcommand{\d}{\ensuremath{\mathrm{d}}}
\setbeamertemplate{caption}{\raggedright\insertcaption\par}

\renewcommand*\familydefault{\sfdefault}
\itemsep2pt
\usepackage[utf8]{inputenc}
\usepackage[T1]{fontenc}
\usepackage{graphicx}
\usepackage{longtable}
\usepackage{wrapfig}
\usepackage{rotating}
\usepackage[normalem]{ulem}
\usepackage{amsmath}
\usepackage{amssymb}
\usepackage{capt-of}
\usepackage{hyperref}
\usetheme{default}
\author{Hjalte Søberg Mikkelsen}
\date{23-04-2026}
\title{The Super Learner}
\begin{document}

\maketitle
\section{Introduction}
\label{sec:orgedca9e1}

\begin{frame}[label={sec:orgfdb0317}]{Problem}
We want to use variables ,\(W\), from the registers to predict the risk of getting diagnosed with type 2 diabetes, \(Y\) (\(Y = 1\) means the person has T2D). To do this we need to find some function \(\psi_0(W) = P(Y = 1)\). \newline
\newline
\pause
How do we find this function?
\end{frame}


\begin{frame}[label={sec:orga357fba}]{Super leaner}
Assume we have the data \(X_i = (W_i,Y_i)\), \(i = 1,...,n\), where \(W_i\) is a p-dimensional set of covariates and \(Y_i\) is the outcome of interest.\pause  Also assume we are given a library of learners \(\mathcal{L} = (\psi_1(W),...,\psi_K(W))\) containing \(K\) learners.\pause The super learner algorithm will then find the function that minimizes the expectation of some loss- or score-function:
\begin{align*}
\psi_0(W) = \arg\min_{\psi\in\mathcal{L}}E[L(X,\psi(W))],
\end{align*}
where the loss function \(L\) is usually the squared loss \(L = (Y-\psi(W))^2\). This is done by using cross validation. 
\end{frame}


\begin{frame}[label={sec:org638d0c2}]{10 fold crossvalidation}
\begin{center}
\includegraphics[width=\linewidth]{/home/hmik0048/PhD-Project/Project_wp2_riskprediction/super_learner_presentation/10foldcrossval.png}
\end{center}
\end{frame}


\begin{frame}[label={sec:org3754756}]{The discrete super learner}
Picks the best performing model contained in the library.

\begin{itemize}
\item Pros:
\begin{itemize}
\item All theoretical results shown by van der Laan and Polley holds \pause
\item Can perform aswell as the oracle estimator \pause
\item Performs well in practise\pause
\item Low runtime\pause
\end{itemize}

\item Cons:
\begin{itemize}
\item Not as fancy as the ensemble learner
\end{itemize}
\end{itemize}
\end{frame}


\begin{frame}[label={sec:orgd80539e}]{The ensemble super learner}
Combines all the learners in the library to create a new "super" learner.

\begin{itemize}
\item Pros:
\begin{itemize}
\item Fancy \pause
\item Usually outperforms the discrete superlearner during training \pause
\end{itemize}

\item Cons:
\begin{itemize}
\item No theoretical results \pause
\item Can never be the best learner \pause
\item Doesnt perform well in practise \pause
\item Crazy bad runtime
\end{itemize}
\end{itemize}
\end{frame}

\begin{frame}[label={sec:org47b1432}]{Evaluate the model}
Normally shown by discrimination- and calibration-plots, and measured by the AUC and the Brier-score, but on what data?\pause
\begin{center}
\includegraphics[width=\linewidth]{/home/hmik0048/PhD-Project/Project_wp2_riskprediction/super_learner_presentation/crossval.png}
\end{center}
\end{frame}



\begin{frame}[label={sec:org53441db}]{Performance measures: AUC}
Measures discrimination, i.e. how good is the model to discriminate between a person who gets diabetes and a person who doesnt. The AUC for a model is the probability of the model, given a random person who gets T2D and a random person who doesnt, to give the T2D person a higher risk. 
\begin{center}
\includegraphics[width=\linewidth]{/home/hmik0048/PhD-Project/Project_wp2_riskprediction/super_learner_presentation/AUC_threshhold.png}
\end{center}
\pause
\newline
Alot of articles use this as a standalone measure (bad). Show examples from google
\newline
\pause
Go to \url{https://developers.google.com/machine-learning/crash-course/classification/roc-and-auc} for an interactive display of the AUC
\end{frame}

\begin{frame}[label={sec:org6f07c72}]{AUC plots}
\begin{center}
\includegraphics[width=100]{/home/hmik0048/PhD-Project/Project_wp2_riskprediction/super_learner_presentation/AUC_worst.png}
\end{center}

\begin{center}
\includegraphics[width=100]{/home/hmik0048/PhD-Project/Project_wp2_riskprediction/super_learner_presentation/AUC_perfect.png}
\end{center}
\end{frame}


\begin{frame}[label={sec:org4bb029c}]{AUC plots}
\begin{center}
\includegraphics[width=\linewidth]{/home/hmik0048/PhD-Project/Project_wp2_riskprediction/super_learner_presentation/AUC_middle.png}
\end{center}
\end{frame}


\begin{frame}[label={sec:org87eb3a8}]{AUC plots}
\begin{center}
\includegraphics[width=100]{/home/hmik0048/PhD-Project/Project_wp2_riskprediction/super_learner_presentation/AUC_lessthan50.png}
\end{center}
\end{frame}


\begin{frame}[label={sec:org3d6aa11}]{Performance measures: Brier-score}
Measure overall performance of a prediction model, and can be used as a standalone measure as it is a theoretical proper measure.
\begin{align*}
BS = \frac{1}{n}\sum_{i = 1}^n(\psi_0(W_i)-Y_i)^2
\end{align*}
\end{frame}



\begin{frame}[label={sec:org6b48336}]{Calibration plots}
Calibration in the large refers to if the model average predicted probability is the same as the observed average probability of the outcome.
\newline
Calibration in the small refers to if the predicted probabilites match the observed probabilites, i.e. 25\% of the people who get a prediction of 25\% should get the event. 
\end{frame}
\end{document}
